
\textbf{Corollaire 6.2.} \uline{Avec les hypothèses de} 6.1, \uline{soit} @@M0@@ \uline{un sous}-@@M1@@-\uline{anneau engendré par les classes d'un nombre fini de faisceaux}. \uline{Alors la filtration de} @@M2@@-\uline{anneau de} @@M3@@ \uline{est discrète}, \uline{et} @@M4@@ \uline{est nilpotent}. \uline{De plus}, @@M5@@ \uline{est un module de type fini sur} @@M6@@ (V 3.9.1).

Comme @@M7@@ est quasi-compact, on peut trouver un recouvrement fini de @@M8@@ par des ensembles ouverts disjoints, tel que sur chacun des ouverts de ce recouvrement les faisceaux dont les classes engendrent @@M9@@ soient tous de rang constant. L'argument standard permet alors de se ramener au cas où ces faisceaux sont de rang constant sur @@M10@@; @@M11@@ est alors augmenté vers @@M12@@.

Soient @@M13@@, @@M14@@, les faisceaux dont les classes engendrent @@M15@@, @@M16@@ leurs classes, @@M17@@. On peut écrire @@M18@@ comme quotient du @@M19@@-anneau engendré par les @@M20@@ soumis aux relations @@M21@@ pour @@M22@@ c'est-à-dire @@M23@@ ; par suite, @@M24@@ est noethérien. Mais d'après la proposition 6.1, les @@M25@@, @@M26@@ sont nilpotents. Donc l'idéal d'augmentation de @@M27@@ est nilpotent; par ailleurs, il n'est autre que @@M28@@.

Tout élément de @@M29@@ s'écrit comme polynome sans composantes de poids @@M30@@ par rapport aux @@M31@@, @@M32@@. Soit @@M33@@ le plus grand entier tel qu'il existe @@M34@@ tel que @@M35@@; soit @@M36@@ un entier tel que les puissances @@M37@@-ièmes des @@M38@@, @@M39@@ soient toutes nulles. Pour @@M40@@, on voit tout de suite que @@M41@@; la filtration de @@M42@@ est donc discrète.

La filtration de @@M43@@ étant discrète, il n'y a qu'un nombre fini de monômes non nuls par rapport aux @@M44@@, @@M45@@. Ces monômes forment un système de générateurs de @@M46@@ comme @@M47@@-module, ce qui achève la démonstration.

